% Generated draft: requires mathematical and source review.
\documentclass[11pt]{article}
\usepackage[margin=25mm]{geometry}
\usepackage{fontspec}
\setmainfont{Latin Modern Roman}
\usepackage{amsmath,amssymb,mathtools,unicode-math}
\setmathfont{Latin Modern Math}
\usepackage{enumitem,needspace,xurl,hyperref}
\setlength{\parindent}{0pt}
\setlength{\parskip}{5pt}
\setlength{\emergencystretch}{3em}
\newcommand{\N}{\mathbb{N}}
\newcommand{\Z}{\mathbb{Z}}
\newcommand{\Q}{\mathbb{Q}}
\newcommand{\R}{\mathbb{R}}
\newcommand{\C}{\mathbb{C}}
\newcommand{\F}{\mathbb{F}}
\newcommand{\A}{\mathbb{A}}
\newcommand{\PP}{\mathbb{P}}
\newcommand{\E}{\mathbb{E}}
\newcommand{\eps}{\varepsilon}
\newcommand{\dd}{\,\mathrm{d}}
\begin{document}
\section*{UnsolvedMath ID: 20003329}
\textbf{Source problem code:} catalogue entry
\subsection{Definitions and mathematical statement}
Let $X=\Gamma\backslash\mathbb H^2$, where $\Gamma$ is a noncocompact arithmetic lattice, and let $V$ be normalized hyperbolic area, $V(X)=1$. The surrounding quantum-chaos setup, recovered from companion records, uses $L^2(V)$-normalized eigenfunctions $\varphi_i$ with Laplace eigenvalues tending to infinity, jointly diagonalizing the relevant invariant differential and Hecke operators, and probabilities
\[
 \dd\mu_i=|\varphi_i|^2\,\dd V.
\]
Problem 40 of E. Lindenstrauss asks for the projective convergence
\[
 \frac{\int_X f\,\dd\mu_i}{\int_X g\,\dd\mu_i}
 \longrightarrow
 \frac{\int_X f\,\dd V}{\int_X g\,\dd V}
 \qquad(f,g\in C_c(X)).
\]
The damaged displayed extraction is a ratio, not a product. The necessary convention is $V(g)\ne0$ and $\mu_i(g)\ne0$ for all sufficiently large $i$; an unrestricted statement for every $g$, including zero, is undefined. Here and below $\rho(f)=\int f\,\dd\rho$. The source's final ``The analog of'' is incomplete.

The supplied aid gives a zero-safe formulation valid for probabilities on a locally compact second-countable space. Fix $h\in C_c(X)$ with $h\ge0$, $V(h)>0$, and $\mu_i(h)>0$ eventually. All admissible ratios above are equivalent to
\[
 \frac{\mu_i}{\mu_i(h)}\longrightarrow\frac{V}{V(h)}
 \quad\text{vaguely, that is, against every }f\in C_c(X).
\]
This anchor-normalized convergence forces every ordinary vague cluster measure to be $cV$ for some $0\le c\le1$. Conversely, the condition that every vague cluster point lies in this ray gives the ratios if $\liminf_i\mu_i(h)>0$, excluding the zero-cluster obstruction. Anchor-normalized convergence together with either $\mu_i(h)\to V(h)$ or tightness of $(\mu_i)$ gives full probability convergence $\mu_i\to V$.

The reported quantitative ratio propagation is explicit. Put
\[
 A_i=\frac{\mu_i(f)}{\mu_i(h)},\quad B_i=\frac{\mu_i(g)}{\mu_i(h)},\quad
 a=\frac{V(f)}{V(h)},\quad b=\frac{V(g)}{V(h)}\ne0.
\]
If $|A_i-a|\le\varepsilon_f$ and $|B_i-b|\le\varepsilon_g\le|b|/2$, then, also for signed or complex tests,
\[
 \left|\frac{A_i}{B_i}-\frac ab\right|
 \le\frac{2\varepsilon_f}{|b|}
      +\frac{2|a|\varepsilon_g}{|b|^2}.
\]
The aid emphasizes that ratios do not recover the missing mass scale. For points $x_i$ leaving every compact set and fixed $0<c<1$, probabilities $cV+(1-c)\delta_{x_i}$ have the required ratios but vague limit $cV$. Similarly, $\varepsilon_i\rho+(1-\varepsilon_i)\delta_{x_i}$ with $\varepsilon_i\downarrow0$ and a fixed probability $\rho$ not proportional to $V$ have only the zero vague cluster, while their local ratios can reflect $\rho$ rather than $V$. The assertion that all nonzero cluster points lie on the $V$-ray would then be vacuous. These measure-theoretic equivalences and examples do not themselves prove arithmetic QUE or no escape for the eigenfunction sequence.
\subsection{Short English statement}
The intended problem asks whether compact-test ratios for arithmetic quantum measures approach the corresponding volume ratios. One positive reference test makes this precise, but ratios determine only a common scale; recovering total mass still requires an anchor-mass limit or no escape.
\subsection{Sources}
\begin{itemize}
\item E. Lindenstrauss, Problem 40, American Institute of Mathematics, \emph{Emerging applications of measure rigidity}, \url{https://aimath.org/WWN/measrigid/measrigid.pdf}.
\item Companion local records in \texttt{aim-other-notes.json}: ID $20003324$, \texttt{AIM-OTHER-0131}, \texttt{notes[130]}, supplies normalized volume and eigenfunction measures; ID $20003325$, \texttt{AIM-OTHER-0132}, \texttt{notes[131]}, supplies the Hecke-eigenfunction assumption; ID $20003328$, \texttt{AIM-OTHER-0135}, \texttt{notes[134]}, explains the noncompact escape issue and cites E. Lindenstrauss [120], with later no-escape work attributed in its aid to Soundararajan. Full bibliography entries are absent.
\item Local UnsolvedMath registry, permanent ID $20003329$, source code \texttt{catalogue entry}, canonical \texttt{aim-other-notes.json}, \texttt{notes[135]}, dated 14 August 2026. Research attempt 1: reduction, solved at run, proved rigor, low novelty confidence; machine-generated and not peer reviewed. The solved label is supplied metadata, not established by the abstract ratio reduction.
\end{itemize}
\end{document}
